% Page 025 — page imprimée 21
% Protestants et Catholiques face à la Révolution à Meyrueis et dans les Cévennes


{\fontsize{10.05}{11.45}\selectfont

La loi prévoyait « la mort pour les chefs et la clémence pour leurs troupes ». Mais le tribunal
de Mende passa outre, après l’arrestation d’un groupe de 40 paysans de la Malène. Trente-
sept furent exécutés à Florac. Trois d’entre eux avaient-ils réussi à s’échapper ? Une plaque
rappelle leur mort dans l’église de la Malène.

\medskip

\textbf{L’engagement protestant}

\vspace{0.35em}

Il ne fait pas de doute que les protestants de Meyrueis ont accueilli favorablement la
Révolution. Nous ne disposons pas, à ma connaissance, de beaucoup de documents en
donnant la preuve si ce n’est un mémoire de Antoine Avesque (1771-1834) que l’on trouve
dans un registre de son fils, Pierre Calixte Avesque, notaire. Il gérait le domaine de Drigas et
avait affermé également, avec trois coreligionnaires (Bosc, Belon et Avesque), le domaine de
Saubert. En voici les premiers paragraphes transmis par Me. Pottier, notaire à Meyrueis, qui a
rajouté la ponctuation, et corrigé l’orthographe :

\begin{center}
\textbf{Au nom de Dieu soit fait Amen}
\end{center}

« Livre de mémoire où on pourra trouver note des affaires majeures et essentielles que j’ai pu traiter, y puiser
bien des renseignements, certains faits et sur diverses affaires…

\begin{enumerate}[label=\arabic*., leftmargin=2.0em, itemsep=0.55em, topsep=0.25em]
\item
La Révolution commença en l’année 1789. Louis XVI ayant convoqué les Etats Généraux, ils finirent
par se rendre maître du pouvoir souverain et par lui faire porter la tête à l’échafaud ainsi que de
perdre toute la famille royale. Quasi toute la noblesse française émigra aux époques de 1791-1792.
Tous les souverains d’Europe se liguèrent contre la France il y eut des guerres terribles, les Français
restèrent toujours vainqueurs, malgré la faiblesse de nos gouvernants, malgré les guerres civiles qui régnaient
dans plusieurs point de France comme à la Vendée, à Lyon, en Lozère, en Ardèche etc. et malgré la trahison de
plusieurs généraux.

Les prêtres qui n’avaient pas voulu prêter le serment de fidélité à l’Etat furent fortement persécutés. On en fit
périr plusieurs. Par la suite les ministres ( les pasteurs) le furent aussi mais d’une manière moins forte.
Dans ce temps on abattit beaucoup de cloches et de clochers, le peuple se porta à des excès de fureur en brûlant
les châteaux et les maisons de certaines personnes qu’on appelait suspects. Enfin la France était dans le plus
grand désordre, l’Histoire en fournira les détails.

\item
A l’époque, au mois de Juillet 1790 je fus député avec m. Ricard par le district de Meyrueis pour aller à Paris à la
grande fédération du 14 Juillet…

\begin{quote}
…A l’époque du 10 Aout 1792, nous partîmes environ \textbf{quatre vingt sept jeunes gens de Meyrueis ville et
Meyrueis campagne}, tous de bonne volonté pour la défense de notre patrie. Nous formâmes notre bataillon à
Mende. Je fus nommé capitaine. J’ai fait quatre campagnes dans les Alpes… Notre bataillon était le premier de
la Lozère*.

\textit{De ces 87 jeunes conscrits combien ont échappé à la mort ?}
\end{quote}
\end{enumerate}

\vspace{0.15em}

Ainsi la 3° compagnie du 1er.Bataillon de l’armée des Alpes était composée
exclusivement de Meyrueisiens. Elle avait à sa tête trois tout jeunes officiers :

\begin{itemize}[leftmargin=1.8em, itemsep=0.18em, topsep=0.2em]
\item Antoine Avesque, 24 ans, capitaine
\item Pellet, \hfill lieutenant
\item Mezin, 1,69m, blessé d’un coup de feu à l’affaire du 28 Messidor an IV sous Mantoue :
la balle lui a traversé la poitrine. A la bataille de Wagram il reçoit deux coups de feu :
l’un lui traverse la jambe droite l’autre lui emporte le bout du nez. Il reçoit la Légion
d’Honneur pour la conduite distinguée qu’il a tenue à Porto-Vecchio en Corse. Il
meurt, chef de bataillon à l’hôpital d’Anvers après 21 ans de service.
\end{itemize}

On trouve encore deux autres meyrueisiens au 3° bataillon

Jean Cambon, 18 ans, Lieutenant et François Vivens, 24 ans, sous-lieutenant qui se
retire à Meyrueis en l’an IV après quatre ans de service. Il est considéré comme
déserteur.

}

